\section{Feature-Set Comparison}
\label{sec:feature_set_results}

The first comparison examines whether conventional ownership and engineered network features improve prediction beyond the retained market variables. The three nested feature sets are evaluated separately within Ridge and XGBoost. Each specification is estimated on the training sample and assessed on the validation sample using a common configuration within its model class. This isolates the information supplied by the features from differences in model flexibility and hyperparameter settings.

\subsection{Ridge Feature-Set Comparison}
\label{sec:linear_feature_results}

Table~\ref{tab:ridge-feature-set-validation} reports the Ridge results, using MKT for market features, OWN for conventional ownership features, and NTW for engineered network features. Adding conventional ownership variables lowers validation MAE only slightly, and the engineered network variables provide a still smaller additional change. AIC and BIC favor the richer specifications because their small in-sample improvements are measured over a large training sample. They are reported as complementary complexity-adjusted diagnostics, while the retained specification is chosen from validation performance and parsimony. On that basis, the additional variables do not produce a material out-of-sample gain.

\begin{table}[!htbp]
\centering
\footnotesize
\setlength{\tabcolsep}{4pt}
\caption{Ridge validation performance and training-sample complexity diagnostics across feature sets.}
\label{tab:ridge-feature-set-validation}
\begin{tabular}{lrrrrrrr}
\toprule
Specification & $\mathrm{MAE}_{\mathrm{train}}$ & $\mathrm{MAE}_{\mathrm{val}}$ & $\mathrm{RMSE}$ & $R^2$ & $\overline{\rho}$ & AIC & BIC \\
\midrule
\textbf{MKT} & \textbf{0.0869} & \textbf{0.0959} & \textbf{0.1337} & \textbf{0.449} & \textbf{0.747} & \textbf{-169,214} & \textbf{-169,136} \\
MKT + OWN & 0.0869 & 0.0956 & 0.1334 & 0.452 & 0.746 & -169,471 & -169,344 \\
MKT + OWN + NTW & 0.0868 & 0.0956 & 0.1332 & 0.453 & 0.746 & -169,637 & -169,452 \\
\bottomrule
\end{tabular}
\end{table}


\noindent
Given these small differences, the market-only feature set is retained for the linear finalist. This is a parsimony decision: the ownership and network variables produce minor numerical improvements in MAE, but not a meaningful change in ranking performance. Validation errors across realized-drawdown deciles for the three Ridge specifications are reported in Appendix~\ref{app:linear_feature_set_diagnostics}.

\subsection{XGBoost Feature-Set Comparison}
\label{sec:nonlinear_feature_results}

Table~\ref{tab:xgboost-feature-set-validation} repeats the comparison using XGBoost under the same initial configuration. Conventional ownership features again produce a small reduction in validation MAE, while adding engineered network features does not improve on the ownership-augmented specification. The three specifications remain very close in both $R^2$ and rank correlation.

\begin{table}[!htbp]
\centering
\footnotesize
\setlength{\tabcolsep}{4pt}
\caption{XGBoost validation performance across feature sets.}
\label{tab:xgboost-feature-set-validation}
\begin{tabular}{lrrrrrr}
\toprule
Specification & Best round & $\mathrm{MAE}_{\mathrm{train}}$ & $\mathrm{MAE}_{\mathrm{val}}$ & $\mathrm{RMSE}$ & $R^2$ & $\overline{\rho}$ \\
\midrule
\textbf{MKT} & \textbf{176} & \textbf{0.0836} & \textbf{0.0950} & \textbf{0.1317} & \textbf{0.466} & \textbf{0.749} \\
MKT + OWN & 175 & 0.0833 & 0.0943 & 0.1305 & 0.475 & 0.751 \\
MKT + OWN + NTW & 140 & 0.0835 & 0.0946 & 0.1306 & 0.474 & 0.750 \\
\bottomrule
\end{tabular}
\end{table}


\noindent
The ownership variables therefore contain some incremental validation signal, but the improvement is too small to justify retaining the larger feature set for the subsequent model-class comparison. The market-only specification is retained and tuned in the next section. This evidence provides limited support for Hypothesis~H1 and no consistent support for Hypothesis~H2. Validation errors across realized-drawdown deciles for the three XGBoost specifications are reported in Appendix~\ref{app:nonlinear_feature_set_diagnostics}.
